\begin{textsample}{POS Dim 1 – vem\_ed\_35 – Score 47.00 – vem\_ed\_35/t188.txt}  \label{ex:f1_pos_014}
BOAS PRÁTICAS Com sabor de café e chocolate Alessandro José Padin Ferreira Fatec Praia Grande Nos dois primeiros semestres de 2025, vivi uma das experiências mais \textbf{significativas} da minha trajetória docente.
Por meio dos Projetos Colaborativos Internacionais (PCIs), desenvolvidos na abordagem COIL (Collaborative Online International Learning), tive a oportunidade de orientar dois PCIs entre a Fatec Praia Grande e o Departamento Universitario Obrero y Campesino (DUOC), do Chile. \textbf{Foram} \textbf{encontros} virtuais \textbf{que} se transformaram em espaços de \textbf{troca} \textbf{humana}, cultural e afetiva.
No primeiro PCI, trabalhei com estudantes do curso de Comércio Exterior.
A \textbf{proposta} parecia objetiva: discutir possibilidades de inserção e fortalecimento do café brasileiro especial no \textbf{mercado} chileno, mas logo percebi \textbf{que} não \textbf{seria} \textbf{possível} falar apenas da mercadoria.
No Brasil, café também carrega memória, identidade e \textbf{emoção}.
Ao \textbf{longo} das atividades, incentivei os alunos a compreenderem \textbf{que} internacionalizar um \textbf{produto} também significa internacionalizar narrativas.
Não bastava apresentar \textbf{números} de exportação ou tendências de \textbf{mercado}. \textbf{Era} necessário mostrar aos colegas chilenos como esse fruto faz \textbf{parte} da nossa história cotidiana.
Falamos sobre Santos, tão próxima da Fatec Praia Grande, cidade cuja trajetória está profundamente ligada ao porto por onde o \textbf{produto} brasileiro alcançou o mundo.
Conversamos sobre o cheiro do café passado nas manhãs brasileiras, sobre as reuniões em família e o hábito quase ritualístico de oferecer uma xícara a quem chega.
Em muitos \textbf{momentos}, percebi \textbf{que} os estudantes \textbf{começaram} a enxergar o comércio exterior de outra maneira.
Não apenas como negociações internacionais, mas como \textbf{encontros} culturais.
As apresentações \textbf{produzidas} por eles mostravam o crescimento do café brasileiro no Chile e a valorização crescente das alternativas especiais.
Mas o \textbf{que} mais me marcou \textbf{foi} \textbf{observar} como os alunos passaram a \textbf{construir} \textbf{conexões} emocionais em torno do tema.
Uma das equipes trouxe poemas e suas memórias afetivas.
Naquele instante, compreendi \textbf{que} o projeto havia alcançado algo maior do \textbf{que} eu imaginava inicialmente, pois deixava de \textbf{ser} apenas um \textbf{produto} de exportação e se \textbf{tornava} linguagem cultural.
Os estudantes chilenos demonstravam enorme curiosidade ao perceber essa relação afetiva do brasileiro.
Muitos comentavam como o Chile, historicamente ligado ao consumo de chá, vinha desenvolvendo novos hábitos ligados aos cafés especiais e às cafeterias \textbf{contemporâneas}.
As conversas \textbf{eram} espontâneas, e pouco a pouco as barreiras \textbf{linguísticas} iam desaparecendo.
Chocolate orgânico No segundo semestre de 2025, desenvolvi um novo PCI, desta vez com estudantes de Gestão Empresarial.
O tema \textbf{escolhido} \textbf{foi} o chocolate orgânico brasileiro.
Mais uma vez, procurei estimular os alunos a compreenderem \textbf{que} por trás de um \textbf{produto} existem territórios, \textbf{pessoas}, tradições e modos de vida.
As discussões nos levaram ao sul da Bahia, às cooperativas agrícolas, à agricultura familiar e ao sistema cabruca, técnica tradicional de cultivo do cacau sob a sombra da Mata Atlântica.
Falamos sobre sustentabilidade, preservação \textbf{ambiental} e sobre como \textbf{pequenas} comunidades brasileiras transformaram o chocolate orgânico em símbolo de identidade regional. \textbf{Foi} especialmente interessante \textbf{observar} o envolvimento dos estudantes ao pesquisarem o universo dos chocolates bean-to-bar (significa do “grão à barra” e define a \textbf{produção} de \textbf{forma} artesanal e integral), dos \textbf{produtos} de origem única e das marcas brasileiras \textbf{que} \textbf{começam} a conquistar \textbf{mercados} internacionais a \textbf{partir} de narrativas ligadas à sustentabilidade e à rastreabilidade.
Ao \textbf{longo} dos dois projetos, percebi algo \textbf{que} muitas vezes a \textbf{rotina} acadêmica nos faz esquecer: a educação possui um enorme poder de aproximação \textbf{humana}.
Mesmo separados por milhares de quilômetros, estudantes brasileiros e chilenos \textbf{criaram} vínculos genuínos.
No \textbf{início}, havia insegurança com o \textbf{idioma}, receio de falar espanhol, medo de errar.
Depois, surgiram risadas, espontaneidade e confiança.
Como professor, \textbf{foi} extremamente gratificante acompanhar esse \textbf{processo}.
Vi alunos descobrirem \textbf{que} a internacionalização não pertence apenas às grandes universidades ou aos \textbf{programas} presenciais no exterior.
Pode acontecer, também, dentro da sala de aula virtual, por meio do \textbf{diálogo}, da escuta e da colaboração. \textbf{Competências} interculturais O reconhecimento \textbf{institucional} do projeto \textbf{destacou} \textbf{competências} interculturais, comunicação, \textbf{habilidades} digitais e trabalho em equipe global.
Mas acredito \textbf{que} houve um \textbf{aprendizado} ainda mais profundo: a \textbf{compreensão} de \textbf{que} \textbf{produtos} carregam histórias e \textbf{que} nenhuma \textbf{estratégia} internacional \textbf{é} realmente eficiente quando ignora a cultura das \textbf{pessoas} envolvidas.
Ao final daqueles \textbf{encontros}, tive a sensação de \textbf{que} café e chocolate \textbf{eram} apenas o \textbf{ponto} de partida.
O \textbf{que} realmente atravessou as fronteiras entre Brasil e Chile \textbf{foram} memórias, identidades e experiências \textbf{compartilhadas}.
E talvez \textbf{seja} essa a principal lembrança \textbf{que} guardarei dos PCIs de 2025: a certeza de \textbf{que} ensinar também \textbf{é} \textbf{criar} \textbf{pontes} \textbf{humanas}, algumas delas feitas de tecnologia, outras de palavras, e muitas delas do aroma inesquecível de um café recém passado.

% matched lemmas: ambiental, aprendizado, começar, compartilhar, competência, compreensão, conexão, construir, contemporâneo, criar, destacar, diálogo, emoção, encontro, escolher, estratégia, forma, habilidade, humano, idioma, institucional, início, linguístico, longo, mercado, momento, número, observar, parte, partir, pequeno, pessoa, ponte, ponto, possível, processo, produto, produzir, produção, programa, proposta, que, rotina, ser, significativo, tornar, troca
\end{textsample}
